\documentclass{beamer}
\usepackage[utf8]{inputenc}
\usepackage[T1]{fontenc}
\usepackage[spanish]{babel}
\usepackage{graphicx} 


\usetheme{Madrid}

\includegraphics[width=0.4\textwidth]{images/logouct.png}


\title{Delicias Temuco Restaurante}
\author{
Nicolás Cayumán, Gabriel Gutiérrez y Ailyn Melillán\\[4pt]
\small Programación 2\\
Profesor: Guido Mellado\\
Ayudante: Joaquín Cantero\\
Ingeniería Civil en Informática\\
Universidad Católica de Temuco
}
\date{21/11/2025}

\begin{document}

% Portada
\begin{frame}
  \titlepage
\end{frame}

\begin{frame}{Mejoras}
    \textbf{Respecto a la anterior evaluación, realizamos las siguientes mejoras:}
    \begin{itemize}
        \item Migración total a ORM (SQLAlchemy).
        \item Implementación de CRUD modulares por entidad.
        \item Creación de DTOs para desacoplar UI y ORM.
        \item Nuevas pestañas en la interfaz (9 módulos nuevos).
        \item Uso de programación funcional (lambda, map, filter, reduce).
        \item Control real de stock mediante recetas (N:N).
        \item Implementación de PDF (carta y boleta).
        \item Gráficos estadísticos con Matplotlib.
        \item Carga masiva CSV y validaciones.
    \end{itemize}

\end{frame}


\begin{frame}{Diagrama de clases}
    \begin{center}
        \includegraphics[width=1\textwidth]{images/diagrama_de_clases.png}
    \end{center}
\end{frame}

\begin{frame}{Diagrama MER}
    \begin{center}
        \includegraphics[width=1\textwidth]{images/modelo_entidad_relacion.png}
    \end{center}
\end{frame}

\begin{frame}{Solución propuesta}
Se implementan cambios variados dentro de la aplicación en Python, que agrega más funciones para la gestión del restaurante:

\begin{itemize}
    
    \item \textbf{Migración hacia ORM:} Se utilizan bases de datos para el manejo y registro de datos.
    \item \textbf{Gestión de clientes:} Se registran clientes para mayor control de ventas.
    \item \textbf{Generación automática de documentos:} boleta y carta en PDF con visualización integrada.
    \item \textbf{Gráficos} se implementa el registro de las compras por dia semana mes y se generan graficos en relacion a esto para el control de ventas.
    \item \textbf{Edición de Menú:}  Se permite desde la misma aplicación editar los menús del restaurante, permitiendo cambios de receta, precios, entre otros.
\end{itemize}
\end{frame}

\section{Solución en Código y Flujo Principal}

\begin{frame}{Flujo Principal del Sistema}
El programa opera mediante una arquitectura basada en:
\begin{itemize}
    \item \textbf{ORM (SQLAlchemy)} para la persistencia.
    \item \textbf{CRUDs modulares} para encapsular reglas de negocio.
    \item \textbf{DTOs} para aislar la interfaz del modelo.
    \item \textbf{CustomTkinter} para la interfaz multipestaña.
\end{itemize}

\bigskip
A continuación se detalla el flujo completo de ejecución.
\end{frame}

\begin{frame}{1. Inicialización de la Aplicación}
\textbf{app.py} es el punto de entrada del sistema:

\begin{enumerate}
    \item Se crean las tablas:
    \begin{itemize}
        \item \texttt{create\_db\_and\_tables()}
    \end{itemize}

    \item Se genera una sesión de base de datos:
    \begin{itemize}
        \item \texttt{self.db\_session = SessionLocal()}
    \end{itemize}

    \item Se inicializan los módulos CRUD:
    \begin{itemize}
        \item \texttt{IngredienteCRUD}, \texttt{MenuCRUD}, \texttt{ClienteCRUD}, \texttt{PedidoCRUD}
    \end{itemize}

    \item Se instancia el \textbf{Pedido (DTO)} que funcionará como carrito temporal.
\end{enumerate}
\end{frame}

\begin{frame}{2. Construcción de la Interfaz Gráfica (UI)}
Al llamar a \texttt{crear\_pestanas()} se generan las vistas:

\begin{itemize}
    \item \textbf{Stock}: registro y actualización de ingredientes.
    \item \textbf{Carga CSV}: carga masiva mediante pandas.
    \item \textbf{Gestión de Menús}: creación, edición y eliminación.
    \item \textbf{Pedido}: selección de menús y validación de stock.
    \item \textbf{Boleta}: generación de PDF (Facade).
    \item \textbf{Carta}: PDF del catálogo de menús.
    \item \textbf{Clientes}: registro y edición.
    \item \textbf{Pedidos}: historial y control.
    \item \textbf{Gráficos}: estadísticas en tiempo real.
\end{itemize}
\end{frame}

\begin{frame}{3. Construcción del Carrito y Lógica de Pedido}
\begin{enumerate}
    \item La UI muestra cada menú como una \textbf{tarjeta} (DTO \texttt{CrearMenu}).
    \item El usuario agrega menús con un botón.
    \item Cada selección:
    \begin{itemize}
        \item Valida stock disponible (\texttt{MenuCRUD.esta\_disponible}).
        \item Aumenta \texttt{cantidad} en el DTO.
        \item Se recalcula el total con \textbf{reduce + lambda}.
    \end{itemize}
    \item Se muestra una lista visual del pedido.
\end{enumerate}

\bigskip
\textbf{Este proceso puede realizarse varias veces antes de confirmar.}
\end{frame}

\begin{frame}{4. Confirmación del Pedido}
Cuando el usuario confirma:

\begin{enumerate}
    \item Se crea un \textbf{PedidoORM} en BD.
    \item Por cada menú del carrito:
    \begin{itemize}
        \item Se inserta un \textbf{DetallePedidoORM}.
        \item Se descuenta el stock según receta (N:N).
    \end{itemize}
    \item Se actualiza el historial de ventas.
\end{enumerate}

\bigskip
\textbf{BoletaFacade} recibe el PedidoORM y genera un PDF con:
\begin{itemize}
    \item Cliente
    \item Fecha
    \item Ítems
    \item Totales
\end{itemize}
\end{frame}

\begin{frame}{5. Generación de PDF y Estadísticas}
\textbf{BoletaFacade.py}:
\begin{itemize}
    \item Construye una estructura de datos para la boleta.
    \item Renderiza el PDF.
    \item Lo muestra con \texttt{CTkPDFViewer}.
\end{itemize}

\bigskip

\textbf{graficos.py}:
\begin{itemize}
    \item Extrae datos desde la BD con ORM.
    \item Genera gráficos de barras y torta en tiempo real.
    \item Integra estos gráficos dentro de la interfaz.
\end{itemize}
\end{frame}


\section{ORM y Programación Funcional}

\begin{frame}{Uso del ORM (Ejemplos)}
    \begin{itemize}
        \item Modelo de datos profesional.
        \item Relaciones entre entidades.
        \item Consultas mediante SQLAlchemy.
        \item Manejo de transacciones (commit/rollback).
        \item Uso de tabla intermedia (MenuIngrediente).
    \end{itemize}
\end{frame}


\begin{frame}{Uso de Funciones Lambda y Funcionales}
    \textbf{Incluido en EV3:}
    \begin{itemize}
        \item \texttt{map} para poblar combobox.
        \item \texttt{filter} para filtrar menús disponibles.
        \item \texttt{reduce} para calcular total del pedido.
        \item \texttt{lambda} como funciones inline.
    \end{itemize}
\end{frame}

\begin{frame}[fragile]{Ejemplos reales}
\textbf{map + lambda:}
\begin{minted}[fontsize=\scriptsize]{python}
etiquetas = list(map(lambda c: f"{c.id} - {c.nombre}",
                     clientes))
\end{minted}

\textbf{reduce:}
\begin{minted}[fontsize=\scriptsize]{python}
total = reduce(lambda t, m: t + m.precio*m.cantidad,
               self.pedido.menus, 0)
\end{minted}
\end{frame}


\begin{frame}{Presentación de la interfaz visual}
    \begin{center}
        \includegraphics[width=1\textwidth]{images/interfaz.png}
    \end{center}
\end{frame}


\end{document}